\chapter{小数}\label{ch:g5:decimals}

十分位和百分位已在\cref{ch:g4:decimals}中出现；
再增加一个数位——千分位，小数体系便完整了。
本章将巩固小数的读法、分解，尤其是比较大小，
并避开那些常见的陷阱。

\section{认识到千分位}

\begin{definition}[小数的数位]\label{def:g5:decimals:places}
\omterm{def:g4:decimals:point}{小数点}之后依次是十分位、百分位、
千分位——每一位的值都是前一位的十分之一：
\[
4.362 = 4 + \frac{3}{10} + \frac{6}{100} + \frac{2}{1000}
= \frac{4\,362}{1000}.
\]
\end{definition}

\begin{omfigure}
\begin{tikzpicture}[scale=0.9]
  \draw (0,0) grid[xstep=2.3, ystep=0.85] (9.2,1.7);
  \node[font=\small] at (1.15,1.275) {个位};
  \node[font=\small] at (3.45,1.275) {十分位};
  \node[font=\small] at (5.75,1.275) {百分位};
  \node[font=\small] at (8.05,1.275) {千分位};
  \node[font=\large, omDef] at (1.15,0.425) {$4$};
  \node[font=\large, omThm] at (3.45,0.425) {$3$};
  \node[font=\large, omThm] at (5.75,0.425) {$6$};
  \node[font=\large, omThm] at (8.05,0.425) {$2$};
  \fill (2.3,0.06) circle (2.5pt);
\end{tikzpicture}

{\small $4.362$ 的数位表：整数部分为蓝色，小数部分为红色，
\omterm{def:g4:decimals:point}{小数点}位于个位与十分位之间。}
\end{omfigure}

\begin{example}[同一个数，多种读法]\label{ex:g5:decimals:readings}
$4.362$ 可读作“四点三六二”，或读作“四又
$362$ 个千分之一”，或分解为 $4 + 0.3 + 0.06 + 0.002$。
末尾添零不改变大小：$4.362 = 4.3620$。但
$4.362 \neq 4.0362$：\emph{中间}的零会把每个数字
推向更小的数位。
\end{example}

\begin{example}[它在数轴上的位置]\label{ex:g5:decimals:line}
确定 $4.362$ 的位置：在 $4$ 与 $5$ 之间；放大后，在 $4.3$ 与
$4.4$ 之间；再次放大，在 $4.36$ 与 $4.37$ 之间，且更靠近 $4.36$。
小数每多一位，就多放大一层。
\end{example}

\begin{omfigure}
\begin{tikzpicture}[scale=1.15]
  \draw[-Stealth] (-0.2,1.6) -- (10.4,1.6);
  \foreach \i in {0,...,10} \draw (\i,1.52) -- (\i,1.68);
  \node[below, font=\small] at (0,1.5) {$4$};
  \node[below, font=\small] at (10,1.5) {$5$};
  \node[below, font=\small] at (3,1.5) {$4.3$};
  \node[below, font=\small] at (4,1.5) {$4.4$};
  \fill[omThm] (3.62,1.6) circle (1.7pt);
  \draw[omExo, thin] (3,1.45) -- (0,0.15);
  \draw[omExo, thin] (4,1.45) -- (10,0.15);
  \draw[-Stealth] (-0.2,0) -- (10.4,0);
  \foreach \i in {0,...,10} \draw (\i,-0.08) -- (\i,0.08);
  \node[below, font=\small] at (0,-0.1) {$4.3$};
  \node[below, font=\small] at (10,-0.1) {$4.4$};
  \node[below, font=\small] at (6,-0.1) {$4.36$};
  \fill[omThm] (6.2,0) circle (1.7pt)
    node[above, font=\small] {$4.362$};
\end{tikzpicture}

{\small 数轴的放大图：从 $4.3$ 到 $4.4$ 这一段，
放大后再分为十个百分之一——而 $4.362$ 位于
$4.36$ 之后不远处。}
\end{omfigure}

\section{比较小数的大小}

\begin{method}[比较两个小数的大小]\label{met:g5:decimals:compare}
\begin{enumerate}
  \item 先比较整数部分：$6.1 > 5.987$。
  \item 整数部分相等时，从\omterm{def:g4:decimals:point}{小数点}起逐位比较小数部分：
        先比十分位，再比百分位，然后
        比千分位；
  \item 用末尾补零的方法，使比较更公平：
        $2.5$ 与 $2.48$ 化为 $2.50$ 与 $2.48$，于是
        $50 > 48$（以百分之一为单位）。
\end{enumerate}
最后再提醒一次这个陷阱：\emph{小数部分的位数更多，
并不代表数更大}——$2.48$ 比 $2.5$ 位数多，却更小。
\end{method}

\begin{example}\label{ex:g5:decimals:order}
将 $7.3$、$7.09$、$7.31$、$7.299$ 排序，补足到三位小数：
$7.300$、$7.090$、$7.310$、$7.299$。于是
\[
7.09 < 7.299 < 7.3 < 7.31 .
\]
注意 $7.299 < 7.3$，尽管 $299$ 看起来很大：
$299$ 个千分之一，对比 $300$ 个千分之一。
\end{example}

\begin{example}[在两个小数之间插入新数]\label{ex:g5:decimals:between}
$5.7$ 与 $5.8$ 之间是否有数？有的，而且很多：$5.75$、
$5.71$、$5.799$\,\dots\ 在 $5.79$ 与 $5.8$ 之间呢？同样很多：
例如 $5.795$。任意两个不同的小数之间，
\emph{总是}能再插入一个数——不存在“紧接着的下一个”
小数。
\end{example}

\section{小数的四舍五入}

\begin{definition}[四舍五入]\label{def:g5:decimals:round}
要\omterm{def:g4:numbers:round}{四舍五入}到个位（或十分位、百分位），只需看其后一位数字：
若该数字大于或等于 $5$，则向前进一位，否则舍去。因此 $4.362$
\omterm{def:g4:numbers:round}{四舍五入}到个位得 $4$，到十分位得 $4.4$，到百分位得 $4.36$。
\end{definition}

\begin{example}[钱款精确到分]\label{ex:g5:decimals:money}
某商品的价格计算结果为 $7.4983$，显示为 $7.50$：现实中的金额
通常\omterm{def:g4:numbers:round}{四舍五入}到百分位（即精确到分）。注意连锁进位：$7.4983 \to
7.50$，其中 $49$ 变成了 $50$——\omterm{def:g4:numbers:round}{四舍五入}可以
连锁影响好几个数位。
\end{example}

\section{练习}

\begin{exercise}[$\star$]\label{exo:g5:decimals:1}
写成小数：$5 + \frac{2}{10} + \frac{7}{1000}$；
\quad $\frac{85}{100}$；\quad $\frac{4\,507}{1000}$；\quad “十二
又千分之九”。
\end{exercise}

\begin{exercise}[$\star$]\label{exo:g5:decimals:2}
在 $27.418$ 中，$4$ 表示几个什么？$8$ 呢？请将这个数
仿照\cref{def:g5:decimals:places}分解。
\end{exercise}

\begin{exercise}[$\star$]\label{exo:g5:decimals:3}
判断下列说法是否正确，并用数位说明理由。
$3.60 = 3.6$; \quad $0.5 = 0.05$; \quad $2.070 = 2.07$; \quad
$1.3 = 1.03$.
\end{exercise}

\begin{exercise}[$\star$]\label{exo:g5:decimals:4}
抄写并填上 $<$、$>$ 或 $=$：
\[
4.8 \;?\; 4.79, \qquad
0.302 \;?\; 0.32, \qquad
6.5 \;?\; 6.500, \qquad
9.09 \;?\; 9.1 .
\]
\end{exercise}

\begin{exercise}[$\star$]\label{exo:g5:decimals:5}
将下列小数从小到大排列：$3.14$；\ $3.4$；\ $3.041$；\
$3.104$; \ $3.41$.
\end{exercise}

\begin{exercise}[$\star$]\label{exo:g5:decimals:6}
每个数位于哪两个相邻的整数之间？又位于哪两个
相邻的一位小数之间？\quad $6.83$；\quad $0.492$；\quad
$12.06$。
\end{exercise}

\begin{exercise}[$\star$]\label{exo:g5:decimals:7}
将 $8.276$ 分别\omterm{def:g4:numbers:round}{四舍五入}到个位、十分位、百分位；将
$3.95$ \omterm{def:g4:numbers:round}{四舍五入}到十分位。
\end{exercise}

\begin{exercise}[$\star$]\label{exo:g5:decimals:8}
写出一个严格介于下列两数之间的数：$2.6$ 与 $2.7$；\quad $0.41$ 与
$0.42$；\quad $5.99$ 与 $6$。
\end{exercise}

\begin{exercise}[$\star$]\label{exo:g5:decimals:9}
四株植物的高度分别是 $0.85$ m、$1.2$ m、$0.9$ m 和
$1.02$ m。将它们按从矮到高排列。
\end{exercise}

\begin{exercise}[$\star\star$]\label{exo:g5:decimals:10}
露西说：“$7.12 > 7.9$，因为 $12 > 9$”。请帮她改正错误，
先比较十分位。
\end{exercise}

\begin{exercise}[$\star\star$]\label{exo:g5:decimals:11}
求所有恰有两位小数、且\omterm{def:g4:numbers:round}{四舍五入}到十分位后等于
$6.4$ 的数。其中最小的数与
最大的数分别是多少？
\end{exercise}
